\documentclass[10pt,a4paper]{amsart}
\usepackage[margin=19mm]{geometry}
\usepackage{amsmath,amssymb,mathtools}
\usepackage[scr=rsfs,cal=euler]{mathalfa}
\usepackage{xfrac}
\usepackage[unicode,hidelinks,bookmarks=false]{hyperref}
\newcommand{\ie}{i.e.\ }
\newcommand{\cf}{cf.\ }
\newcommand{\resp}{respectively\ }
\newcommand{\Subsection}{\subsection}
\providecommand{\nobreakdash}{\nobreak\hbox{-}}
\setlength{\emergencystretch}{2em}
\multlinegap0pt
\usepackage{bm}
\usepackage{bbm}
\usepackage{dsfont}
\usepackage{xspace}
\usepackage{cancel}
\usepackage{mathtools}
\usepackage{amsthm}
\makeatletter
\@ifundefined{proposition}{\newtheorem{proposition}{Proposition}}{}
\makeatother
\DeclareMathOperator{\Cov}{Cov}
\DeclareMathOperator{\Var}{Var}
\title[A Relativity-Based Framework for Statistical Testing Guided ]{A Relativity-Based Framework for Statistical Testing Guided by the Independence of Ancillary Statistics: Methodology and Nonparametric Illustrations}
\author{Albert Vexler, Douglas Landsittel}
\date{Source: arXiv:2507.11816v2 (CC BY 4.0)}
\begin{document}
\maketitle

\noindent\emph{Source and licence.} A Relativity-Based Framework for Statistical Testing Guided by the Independence of Ancillary Statistics: Methodology and Nonparametric Illustrations. Version arXiv:2507.11816v2 (CC BY 4.0), distributed under the Creative Commons Attribution 4.0 International licence, as recorded in the arXiv licence metadata for this exact version and shown on the abstract page https://arxiv.org/abs/2507.11816v2. Full rights notice: licenses/arxiv-2507.11816-CC-BY-4.0.txt.\par\smallskip

\noindent\textit{Editorial scope.} Counted unit: the Proposition whose source label is \texttt{thm:ratio} in the article (source0.tex lines 339--375, source label \texttt{thm:ratio}), delivered together with the article's own complete original proof of it (source0.tex lines 377--438, one \texttt{proof} environment of 62 lines). No mathematics is written, repaired or re-derived: statement and proof are the article's own text line for line. Required context retained uncounted: none. Every definition, display and label of those lines is retained unchanged. Omitted from this package and not redistributed: 1--338, 376--376, 439--451. Nothing inside a retained excerpt is omitted or altered: every retained body is delivered exactly as the article's lines are written, so no omission receipt of any kind accompanies this package. Citations: the retained excerpts carry no citation command at all, so no bibliography entry of the article is transcribed and no reference is redistributed. Reference adaptation: none was needed and none was made; every reference inside the delivered unit resolves inside it, so no unresolved reference or citation is produced. Printed slips kept literal: no printed word, symbol, hypothesis or number of the counted statement or of its proof was altered. Layout and class: the article's own class is \texttt{article} with packages including amsmath, amssymb, amsthm, bm, graphicx, natbib, geometry, xr; the wrapper is the lane's pinned \texttt{amsart} editorial wrapper at 10pt on A4, which restates the theorem-like environments the retained text needs and adds the article's own macro definitions that the retained text needs (\texttt{\textbackslash Cov}, \texttt{\textbackslash Var}), with extra wrapper packages (bm, bbm, dsfont, xspace, cancel, mathtools). The delivered numbering is the unit's own. Assets: no retained span contains a figure environment, a graphics-inclusion command, a PDF-page inclusion, a table or a TikZ picture; no image file is copied into this package, the package declares no asset dependency and no image file is redistributed with it. Licence: version arXiv:2507.11816v2 of this article is distributed under the Creative Commons Attribution 4.0 International licence, as recorded in the arXiv licence metadata for this exact version and shown on the abstract page https://arxiv.org/abs/2507.11816v2. A full rights notice is delivered at licenses/arxiv-2507.11816-CC-BY-4.0.txt. Characters: every retained excerpt is pure ASCII, so the delivered engine, XeLaTeX with its default Latin Modern text font, reports no missing character.
\begin{proposition} [The asymptotic approximation of $R$]
	\label{thm:ratio}
	Let $\mathbf{X}_{1},\ldots,\mathbf{X}_{n}$ be i.i.d.\ $p$-dimensional random vectors with 
	$\mathrm{E}(\mathbf{X}_{i})=\mu$, $\mathrm{Cov}(\mathbf{X}_{i})=\Sigma$, and finite fourth moments. 
	Assume the elliptical fourth-moment identity
	\begin{equation}
		\mathrm{E}\!\left\{ (Y^\top B Y)(Y^\top C Y) \right\}
		= 2\,\mathrm{tr}(B\Sigma C\Sigma) + (\theta-1)\,\mathrm{tr}(B\Sigma)\,\mathrm{tr}(C\Sigma),
		\qquad Y=\mathbf{X}_1-\mu,
		\label{eq:elliptical-identity}
	\end{equation}
	for symmetric $p\times p$ matrices $B,C$, where $\theta\ge 1$ is the elliptical kurtosis factor. 
	Let
	\[
	T^{2}=n\,\bar{\mathbf{X}}^{\top} S^{-1}\bar{\mathbf{X}}, 
	\qquad 
	T_{1}^{2}=\frac{n p\,\bar{\mathbf{X}}^{\top}\bar{\mathbf{X}}}{\mathrm{tr}(S)},
	\qquad
	\psi=\frac{p\,\mathrm{tr}(\Sigma^2)}{\{\mathrm{tr}(\Sigma)\}^{2}}\ge 1,
	\]
	and denote $W=VV^\top$, where $V=[\,\mathbf{X}_{1}-\bar{\mathbf{X}},\ldots,\mathbf{X}_{n}-\bar{\mathbf{X}}\,]^{\top}$. 
	Then, as $n\to\infty$,
	\begin{equation}
		\frac{\mathrm{RV}(T_1^2,\mathrm{vec}(W))}{\mathrm{RV}(T^2,\mathrm{vec}(W))}
		=\frac{\big(2\psi+(\theta-1)p\big)^2}{\psi\,\big(2+(\theta-1)p\big)^2} + r_{n,p},
		\label{eq:ratio-main}
	\end{equation}
	where the remainder $r_{n,p}$ satisfies
	\[
	r_{n,p}=
	\begin{cases}
		O(n^{-1}), & \text{if $p$ is fixed}, \\[0.5ex]
		O(n^{-1/2}), & \text{if $p,n\to\infty$ with $p/n\to c\in(0,1)$}.
	\end{cases}
	\]
	The leading term in \eqref{eq:ratio-main} depends only on $\psi$ and $\theta$ and is free of $\mu$.
\end{proposition}

\begin{proof} We outline the main steps leading to the expansion of $R$.
	\textit{Step 1: Linearizations of $T^2$ and $T_1^2$.}  
	Let $S=\Sigma+\Delta$ with $\Delta=O_p(n^{-1/2})$ when $p$ is fixed. Expanding yields
	\begin{align}
		T^2 &= n\,\bar{\mathbf{X}}^{\top}\Sigma^{-1}\bar{\mathbf{X}}
		- n\,\bar{\mathbf{X}}^{\top}\Sigma^{-1}\Delta\Sigma^{-1}\bar{\mathbf{X}} + R_{T}, \label{eq:T2-exp}\\
		T_1^2 &= \frac{n p}{\mathrm{tr}(\Sigma)}\,\bar{\mathbf{X}}^{\top}\bar{\mathbf{X}}
		- \frac{n p}{\{\mathrm{tr}(\Sigma)\}^{2}}\,\bar{\mathbf{X}}^{\top}\bar{\mathbf{X}}\,\mathrm{tr}(\Delta) + R_{T1}, \label{eq:T1-exp}
	\end{align}
	with $R_T,R_{T1}=O_p(n^{-1})$. Thus $T^2$ depends on $\Delta$ only through $L_2(\Delta)=\mathrm{tr}(\Sigma^{-1}\Delta)$, while $T_1^2$ depends on $\Delta$ only through $L_1(\Delta)=\mathrm{tr}(\Delta)$.
	
	\textit{Step 2: Covariances with $\mathrm{vec}(W)$.}  
	Since $S=(n-1)^{-1}V^\top V$ and $W=VV^\top$, covariances between $L_1(\Delta)$ or $L_2(\Delta)$ and $\mathrm{vec}(W)$ can be expressed using fourth moments of $Y=\mathbf{X}_i-\mu$. 
	By \eqref{eq:elliptical-identity},
	\[
	\mathrm{E}\{(Y^\top BY)(Y^\top CY)\}
	=2\,\mathrm{tr}(B\Sigma C\Sigma)+(\theta-1)\,\mathrm{tr}(B\Sigma)\,\mathrm{tr}(C\Sigma).
	\]
	Applying this with $B=C$ equal to $I_p$ and $\Sigma^{-1}$, respectively, gives
	\[
	\Gamma_{I_p} = 2\,\mathrm{tr}(\Sigma^2) + (\theta-1)\{\mathrm{tr}(\Sigma)\}^2, 
	\qquad
	\Gamma_{\Sigma^{-1}} = 2p+(\theta-1)p^2.
	\]
	Hence, up to order $O(n^{-1})$, 
	\[
	\Cov(T_1^2,\mathrm{vec}(W)) \propto \Gamma_{I_p}, 
	\qquad 
	\Cov(T^2,\mathrm{vec}(W)) \propto \Gamma_{\Sigma^{-1}}.
	\]
	
	\textit{Step 3: Variances.}  
	For fixed $p$, standard quadratic-form CLTs imply
	\[
	\Var\!\left(n\,\bar{\mathbf{X}}^{\top}\Sigma^{-1}\bar{\mathbf{X}}\right) = 2p+O(n^{-1}),
	\qquad
	\Var\!\left(\tfrac{n p}{\mathrm{tr}(\Sigma)}\,\bar{\mathbf{X}}^{\top}\bar{\mathbf{X}}\right) 
	= 2p\,\psi+O(n^{-1}).
	\]
	For proportional growth $p/n\to c\in(0,1)$, both variances inflate by a factor of $(1-c)^{-3}$, and fluctuations scale as $O(n^{-1/2})$.
	
	\textit{Step 4: RV ratios.}  
	By the definition of RV,
	\[
	\mathrm{RV}(U,\mathbf{Y})=\frac{\|\Cov(U,\mathbf{Y})\|^2}{\Var(U)\,\Var(\mathbf{Y})}.
	\]
	Since the common factor $\Var(\mathrm{vec}(W))$ cancels in the ratio, combining Steps 2 and 3 gives
	\[
	\frac{\mathrm{RV}(T_1^2,\mathrm{vec}(W))}{\mathrm{RV}(T^2,\mathrm{vec}(W))}
	=\frac{\{\Gamma_{I_p}\}^2/\Var(T_1^2)}{\{\Gamma_{\Sigma^{-1}}\}^2/\Var(T^2)}+r_{n,p}.
	\]
	Substituting $\Gamma_{I_p},\Gamma_{\Sigma^{-1}}$ and the variances yields
	\[
	\frac{\mathrm{RV}(T_1^2,\mathrm{vec}(W))}{\mathrm{RV}(T^2,\mathrm{vec}(W))}
	=\frac{\big(2\psi+(\theta-1)p\big)^2}{\psi\,\big(2+(\theta-1)p\big)^2}+r_{n,p},
	\]
	with $r_{n,p}=O(n^{-1})$ for fixed $p$, and $r_{n,p}=O(n^{-1/2})$ for proportional growth $p/n\to c\in(0,1)$.
	
	\textit{Step 5: Independence of $\mu$.}  
	The leading term in \eqref{eq:ratio-main} does not involve $\mu$, since $\mu$ affects only the mean of $\bar{\mathbf{X}}$ while the quadratic-form variances and fourth-moment identities depend solely on centered quantities.
	
\end{proof}

\end{document}
